% GENERATED FILE. Edit canonical lessons, then run npm run generate:books.
% canonical-source-hash: fnv1a64:bbb51a7830bfaa52
% canonical-lessons: TE-C421-apayam, TE-C421-mundujagratta, TE-C421-jagrattapadu, TE-C421-nivarincu, TE-C421-sandigdham, TE-R421-review-414-421

\chapter{Risk and Caution}
\label{ch:te-risk-and-caution}
\hlchaptermodality{421}

\begin{chapteropening}
\textbf{By the end of this chapter:} \emph{I can talk about dangers, precautions and preventing harm.}

Complete the last lesson of chapter 421: I can talk about dangers, precautions and preventing harm.
\end{chapteropening}

\section[\texorpdfstring{\te{అపాయం} (apāyaṁ)}{అపాయం (apāyaṁ)}]{\te{అపాయం} (apāyaṁ) — a danger; a hazard}

\label{lesson:TE-C421-apayam}



\begin{quote}
Before the new one: say the Telugu for optional, then the Telugu for to let go of.
\end{quote}

\subsection*{You'll want to know: \te{అపాయం}}
\textbf{\te{అపాయం}} — \emph{apāyaṁ} — \textquotedblleft{}a danger; a hazard\textquotedblright{}.

\begin{quote}
\te{వరదల} \te{సమయంలో} \te{నదిలోకి} \te{వెళ్ళడం} \te{అపాయం}.
\end{quote}

\begin{quote}
\emph{varadala samayaṁlō nadilōki veḷḷaḍaṁ apāyaṁ.}
\end{quote}

\begin{quote}
Going into the river during floods is dangerous.
\end{quote}

\te{ప్రమాదం}, which you know, is an accident or a danger; \te{అపాయం} is the danger itself.

\begin{cousinweb}
\textbf{\te{అపాయం}} (\emph{apāyaṁ}) is borrowed from Sanskrit \emph{apāya}, \textquotedblleft{}departure, death, misfortune\textquotedblright{} (CDIAL 481).
\end{cousinweb}

\begin{grammarlens}[title={What you've built}]
One more word for saying what happens if.
\end{grammarlens}

\subsection*{Your turn}
\begin{itemize}
\raggedright
  \item \emph{Say it:} \emph{apāyaṁ}
  \item \emph{Say it:} \emph{apāyaṁ}, once more
  \item \emph{Say it:} \emph{varadala samayaṁlō nadilōki veḷḷaḍaṁ apāyaṁ.} — Going into the river during floods is dangerous.
  \item \emph{Recall:} say the Telugu for optional, then the Telugu for to let go of, then say \emph{apāyaṁ} again
\end{itemize}

\begin{tcolorbox}[breakable,colback=teal!4,colframe=teal!35!black,title={Before you move on}]
What does \te{అపాయం} mean? (\textquotedblleft{}A danger; a hazard\textquotedblright{}.) Say it once more.
\end{tcolorbox}
\section[\texorpdfstring{\te{ముందుజాగ్రత్త} (mundujāgratta)}{ముందుజాగ్రత్త (mundujāgratta)}]{\te{ముందుజాగ్రత్త} (mundujāgratta) — a precaution}

\label{lesson:TE-C421-mundujagratta}



\begin{quote}
Before the new one: say the Telugu for to let go of, then the Telugu for a danger.
\end{quote}

\subsection*{You'll want to know: \te{ముందుజాగ్రత్త}}
\textbf{\te{ముందుజాగ్రత్త}} — \emph{mundujāgratta} — \textquotedblleft{}a precaution\textquotedblright{}.

\begin{quote}
\te{ముందుజాగ్రత్తగా} \te{నేను} \te{గొడుగు} \te{తీసుకెళ్ళాను}.
\end{quote}

\begin{quote}
\emph{mundujāgrattagā nēnu goḍugu tīsukeḷḷānu.}
\end{quote}

\begin{quote}
As a precaution I took an umbrella.
\end{quote}

\te{ముందు}, before, + \te{జాగ్రత్త}, care, which you know.

\begin{cousinweb}
\textbf{\te{ముందుజాగ్రత్త}} (\emph{mundujāgratta}) is built from native \emph{mundu}, \textquotedblleft{}in front, before\textquotedblright{} (DEDR 5020), and \te{జాగ్రత్త}, \textquotedblleft{}care\textquotedblright{}, from Sanskrit \emph{jāgrat}, \textquotedblleft{}waking, watchful\textquotedblright{} (CDIAL 5174).
\end{cousinweb}

\begin{grammarlens}[title={What you've built}]
One more word for saying what happens if.
\end{grammarlens}

\subsection*{Your turn}
\begin{itemize}
\raggedright
  \item \emph{Say it:} \emph{mundujāgratta}
  \item \emph{Say it:} \emph{mundujāgratta}, once more
  \item \emph{Say it:} \emph{mundujāgrattagā nēnu goḍugu tīsukeḷḷānu.} — As a precaution I took an umbrella.
  \item \emph{Recall:} say the Telugu for to let go of, then the Telugu for a danger, then say \emph{mundujāgratta} again
\end{itemize}

\begin{tcolorbox}[breakable,colback=teal!4,colframe=teal!35!black,title={Before you move on}]
What does \te{ముందుజాగ్రత్త} mean? (\textquotedblleft{}A precaution\textquotedblright{}.) Say it once more.
\end{tcolorbox}
\section[\texorpdfstring{\te{జాగ్రత్తపడు} (jāgrattapaḍu)}{జాగ్రత్తపడు (jāgrattapaḍu)}]{\te{జాగ్రత్తపడు} (jāgrattapaḍu) — to take care; to be on guard}

\label{lesson:TE-C421-jagrattapadu}



\begin{quote}
Before the new one: say the Telugu for a danger, then the Telugu for a precaution.
\end{quote}

\subsection*{You'll want to know: \te{జాగ్రత్తపడు}}
\textbf{\te{జాగ్రత్తపడు}} — \emph{jāgrattapaḍu} — \textquotedblleft{}to take care; to be on guard\textquotedblright{}.

\begin{quote}
\te{డబ్బు} \te{పోకుండా} \te{నేను} \te{జాగ్రత్తపడ్డాను}.
\end{quote}

\begin{quote}
\emph{ḍabbu pōkuṇḍā nēnu jāgrattapaḍḍānu.}
\end{quote}

\begin{quote}
I took care not to lose the money.
\end{quote}

Past (I …): \te{జాగ్రత్తపడ్డాను}, \emph{jāgrattapaḍḍānu}.

What you guard against takes -\te{కుండా}: \te{పోకుండా}, so that it is not lost.

\begin{cousinweb}
\textbf{\te{జాగ్రత్తపడు}} (\emph{jāgrattapaḍu}) is built from \te{జాగ్రత్త}, \textquotedblleft{}care\textquotedblright{}, from Sanskrit \emph{jāgrat}, \textquotedblleft{}waking\textquotedblright{} (CDIAL 5174), and the native verb \emph{paḍu}, \textquotedblleft{}to feel, to undergo\textquotedblright{} (DEDR 3853).
\end{cousinweb}

\begin{grammarlens}[title={What you've built}]
One more word for saying what happens if.
\end{grammarlens}

\subsection*{Your turn}
\begin{itemize}
\raggedright
  \item \emph{Say it:} \emph{jāgrattapaḍu}
  \item \emph{Say it:} \emph{jāgrattapaḍu}, once more
  \item \emph{Say it:} \emph{ḍabbu pōkuṇḍā nēnu jāgrattapaḍḍānu.} — I took care not to lose the money.
  \item \emph{Recall:} say the Telugu for a danger, then the Telugu for a precaution, then say \emph{jāgrattapaḍu} again
\end{itemize}

\begin{tcolorbox}[breakable,colback=teal!4,colframe=teal!35!black,title={Before you move on}]
What does \te{జాగ్రత్తపడు} mean? (\textquotedblleft{}To take care; to be on guard\textquotedblright{}.) Say it once more.
\end{tcolorbox}
\section[\texorpdfstring{\te{నివారించు} (nivāriñcu)}{నివారించు (nivāriñcu)}]{\te{నివారించు} (nivāriñcu) — to prevent; to avert}

\label{lesson:TE-C421-nivarincu}



\begin{quote}
Before the new one: say the Telugu for a precaution, then the Telugu for to take care.
\end{quote}

\subsection*{You'll want to know: \te{నివారించు}}
\textbf{\te{నివారించు}} — \emph{nivāriñcu} — \textquotedblleft{}to prevent; to avert\textquotedblright{}.

\begin{quote}
\te{పరిశుభ్రత} \te{వల్ల} \te{చాలా} \te{రోగాలను} \te{నివారించవచ్చు}.
\end{quote}

\begin{quote}
\emph{pariśubhrata valla cālā rōgālanu nivāriñcavaccu.}
\end{quote}

\begin{quote}
Many diseases can be prevented by cleanliness.
\end{quote}

Past (I …): \te{నివారించాను}, \emph{nivāriñcānu}.

The noun is \te{నివారణ}, prevention.

\begin{cousinweb}
\textbf{\te{నివారించు}} (\emph{nivāriñcu}) is built on Sanskrit \emph{nivāraṇa}, \textquotedblleft{}keeping back, preventing\textquotedblright{} (CDIAL 7418), with the Telugu verb ending -\te{ఇంచు}.
\end{cousinweb}

\begin{grammarlens}[title={What you've built}]
One more word for saying what happens if.
\end{grammarlens}

\subsection*{Your turn}
\begin{itemize}
\raggedright
  \item \emph{Say it:} \emph{nivāriñcu}
  \item \emph{Say it:} \emph{nivāriñcu}, once more
  \item \emph{Say it:} \emph{pariśubhrata valla cālā rōgālanu nivāriñcavaccu.} — Many diseases can be prevented by cleanliness.
  \item \emph{Recall:} say the Telugu for a precaution, then the Telugu for to take care, then say \emph{nivāriñcu} again
\end{itemize}

\begin{tcolorbox}[breakable,colback=teal!4,colframe=teal!35!black,title={Before you move on}]
What does \te{నివారించు} mean? (\textquotedblleft{}To prevent; to avert\textquotedblright{}.) Say it once more.
\end{tcolorbox}
\section[\texorpdfstring{\te{సందిగ్ధం} (sandigdhaṁ)}{సందిగ్ధం (sandigdhaṁ)}]{\te{సందిగ్ధం} (sandigdhaṁ) — a dilemma; indecision}

\label{lesson:TE-C421-sandigdham}



\begin{quote}
Before the new one: say the Telugu for to take care, then the Telugu for to prevent.
\end{quote}

\subsection*{You'll want to know: \te{సందిగ్ధం}}
\textbf{\te{సందిగ్ధం}} — \emph{sandigdhaṁ} — \textquotedblleft{}a dilemma; indecision\textquotedblright{}.

\begin{quote}
\te{ఏ} \te{ఉద్యోగం} \te{ఎంచుకోవాలో} \te{తెలియక} \te{నేను} \te{సందిగ్ధంలో} \te{ఉన్నాను}.
\end{quote}

\begin{quote}
\emph{ē udyōgaṁ eñcukōvālō teliyaka nēnu sandigdhaṁlō unnānu.}
\end{quote}

\begin{quote}
Not knowing which job to choose, I am in a dilemma.
\end{quote}

\te{సందిగ్ధంలో} \te{ఉండు}, to be in two minds.

\begin{cousinweb}
\textbf{\te{సందిగ్ధం}} (\emph{sandigdhaṁ}) is borrowed from Sanskrit \emph{saṁdigdha}, \textquotedblleft{}confused; questionable, doubtful\textquotedblright{} (Monier-Williams), of the same family as \te{సందేహం}.
\end{cousinweb}

\begin{grammarlens}[title={What you've built}]
One more word for saying what happens if.
\end{grammarlens}

\subsection*{Your turn}
\begin{itemize}
\raggedright
  \item \emph{Say it:} \emph{sandigdhaṁ}
  \item \emph{Say it:} \emph{sandigdhaṁ}, once more
  \item \emph{Say it:} \emph{ē udyōgaṁ eñcukōvālō teliyaka nēnu sandigdhaṁlō unnānu.} — Not knowing which job to choose, I am in a dilemma.
  \item \emph{Recall:} say the Telugu for to take care, then the Telugu for to prevent, then say \emph{sandigdhaṁ} again
\end{itemize}

\begin{tcolorbox}[breakable,colback=teal!4,colframe=teal!35!black,title={Before you move on}]
What does \te{సందిగ్ధం} mean? (\textquotedblleft{}A dilemma; indecision\textquotedblright{}.) Say it once more.
\end{tcolorbox}
\section[(dialogue)]{Words from chapters 414–421 — a review}

\label{lesson:TE-R421-review-414-421}



\begin{quote}
Say the Telugu for a professor and a lecture.
\end{quote}

\subsection*{Your turn}
Say these aloud:

\begin{itemize}
\raggedright
  \item a professor, a lecture, a graduate, a scholarship, to question
  \item education, knowledge, study, to debate, to instruct
  \item a textbook, an essay, a scholarly book, a syllabus topic, to review
  \item a stipulation, eligibility, to permit, except for, to apply
  \item possible, impossible, probability, unavoidable, if possible
  \item if not, depending on, an exemption, in any case, to occur
  \item a choice, an alternative, to take into account, optional, to let go of
  \item a danger, a precaution, to take care, to prevent, a dilemma
\end{itemize}

\begin{tcolorbox}[breakable,colback=teal!4,colframe=teal!35!black,title={Before you move on}]
A professor? (\textbf{\te{ప్రొఫెసరు}}, \emph{prophesaru}.) A lecture? (\textbf{\te{ఉపన్యాసం}}, \emph{upanyāsaṁ}.) A graduate? (\textbf{\te{పట్టభద్రుడు}}, \emph{paṭṭabhadruḍu}.) A scholarship? (\textbf{\te{ఉపకారవేతనం}}, \emph{upakāravētanaṁ}.) To question? (\textbf{\te{ప్రశ్నించు}}, \emph{praśniñcu}.) Education? (\textbf{\te{విద్య}}, \emph{vidya}.) Knowledge? (\textbf{\te{జ్ఞానం}}, \emph{jñānaṁ}.) Study? (\textbf{\te{అధ్యయనం}}, \emph{adhyayanaṁ}.) To debate? (\textbf{\te{తర్కించు}}, \emph{tarkiñcu}.) To instruct? (\textbf{\te{బోధించు}}, \emph{bōdhiñcu}.) A textbook? (\textbf{\te{పాఠ్యపుస్తకం}}, \emph{pāṭhyapustakaṁ}.) An essay? (\textbf{\te{వ్యాసం}}, \emph{vyāsaṁ}.) A scholarly book? (\textbf{\te{గ్రంథం}}, \emph{granthaṁ}.) A syllabus topic? (\textbf{\te{పాఠ్యాంశం}}, \emph{pāṭhyāṁśaṁ}.) To review? (\textbf{\te{సమీక్షించు}}, \emph{samīkṣiñcu}.) A stipulation? (\textbf{\te{నిబంధన}}, \emph{nibandhana}.) Eligibility? (\textbf{\te{అర్హత}}, \emph{arhata}.) To permit? (\textbf{\te{అనుమతించు}}, \emph{anumatiñcu}.) Except for? (\textbf{\te{మినహా}}, \emph{minahā}.) To apply? (\textbf{\te{వర్తించు}}, \emph{vartiñcu}.) Possible? (\textbf{\te{సాధ్యం}}, \emph{sādhyaṁ}.) Impossible? (\textbf{\te{అసాధ్యం}}, \emph{asādhyaṁ}.) Probability? (\textbf{\te{సంభావ్యత}}, \emph{sambhāvyata}.) Unavoidable? (\textbf{\te{అనివార్యం}}, \emph{anivāryaṁ}.) If possible? (\textbf{\te{వీలైతే}}, \emph{vīlaitē}.) If not? (\textbf{\te{కాకపోతే}}, \emph{kākapōtē}.) Depending on? (\textbf{\te{బట్టి}}, \emph{baṭṭi}.) An exemption? (\textbf{\te{మినహాయింపు}}, \emph{minahāyimpu}.) In any case? (\textbf{\te{ఏమైనప్పటికీ}}, \emph{ēmainappaṭikī}.) To occur? (\textbf{\te{సంభవించు}}, \emph{sambhaviñcu}.) A choice? (\textbf{\te{ఎంపిక}}, \emph{empika}.) An alternative? (\textbf{\te{ప్రత్యామ్నాయం}}, \emph{pratyāmnāyaṁ}.) To take into account? (\textbf{\te{పరిగణించు}}, \emph{parigaṇiñcu}.) Optional? (\textbf{\te{ఐచ్ఛికం}}, \emph{aicchikaṁ}.) To let go of? (\textbf{\te{వదులుకో}}, \emph{vadulukō}.) A danger? (\textbf{\te{అపాయం}}, \emph{apāyaṁ}.) A precaution? (\textbf{\te{ముందుజాగ్రత్త}}, \emph{mundujāgratta}.) To take care? (\textbf{\te{జాగ్రత్తపడు}}, \emph{jāgrattapaḍu}.) To prevent? (\textbf{\te{నివారించు}}, \emph{nivāriñcu}.) A dilemma? (\textbf{\te{సందిగ్ధం}}, \emph{sandigdhaṁ}.)
\end{tcolorbox}
